\BourbakiExerciseGroup

\begin{enumerate}
\def\labelenumi{\arabic{enumi})}
\tightlist
\item
  设 \(u(\mathbf{X}) = \sum_{n=0}^{\infty} \alpha_n \mathbf{X}^n\) 是交换域 \(\mathbf{K}\) 上的一个形式幂级数。
\end{enumerate}

\begin{enumerate}
\def\labelenumi{\alph{enumi})}
\tightlist
\item
  要使 u 成为 \(\mathbf{K}(\mathbf{X})\) 中的有理分式，必要且充分的是：存在由 K 的元素组成的不全为零的有限序列 \((\lambda_{i})_{1\leq i\leq q}\)，以及一个整数 \(d\geq0\)，使得对所有 \(n\geq d\) 有
\end{enumerate}

\[
\lambda_ {1} \alpha_ {n} + \lambda_ {2} \alpha_ {n + 1} + \dots + \lambda_ {q} \alpha_ {n + q - 1} = 0.
\]

\begin{enumerate}
\def\labelenumi{\alph{enumi})}
\setcounter{enumi}{1}
\tightlist
\item
  令
\end{enumerate}

\[
\mathrm{H} _ {n} ^ {(k)} = \left| \begin{array}{c c c c} \alpha_ {n} & \alpha_ {n + 1} & \dots & \alpha_ {n + k - 1} \\ \alpha_ {n + 1} & \alpha_ {n + 2} & \dots & \alpha_ {n + k} \\ \alpha_ {n + 2} & \alpha_ {n + 3} & \dots & \alpha_ {n + k + 1} \\ \vdots & \vdots & \ddots & \vdots \\ \alpha_ {n + k - 1} & \alpha_ {n + k} & \dots & \alpha_ {n + 2 k - 2} \end{array} \right|
\]

（“汉克尔行列式”）证明：如果 \(\mathrm{H}_{d+j}^{(q+1)} = 0\) 且 \(\mathrm{H}_{d+j}^{(q)} \neq 0\) 对每个整数 \(j \geqslant 0\) 成立，则 \(u(\mathbf{X})\) 是一个有理分式（利用 \(a\)）。

\begin{enumerate}
\def\labelenumi{\alph{enumi})}
\setcounter{enumi}{2}
\tightlist
\item
  证明恒等式
\end{enumerate}

\[
\mathrm{H} _ {n} ^ {(k)} \mathrm{H} _ {n + 2} ^ {(k)} - \mathrm{H} _ {n} ^ {(k + 1)} \mathrm{H} _ {n + 2} ^ {(k - 1)} = (\mathrm{H} _ {n + 1} ^ {(k)}) ^ {2}
\]

（参见 III，第 639 页，练习 10）。由此推出，如果 \(\mathrm{H}_{m+j}^{(k+1)} = 0\) 对 \(0 \leqslant j \leqslant r-1\) 成立，则 \(r\) 个行列式 \(\mathrm{H}_{m+j}^{(k)}\) （其中 \(1 \leqslant j \leqslant r\)）要么全为零，要么全不为零。

\(d\)）由 \(b\)）和 \(c\)）推出，要使 \(u(\mathbf{X})\) 是有理分式，必要且充分的是存在两个整数 \(d\) 和 \(q\)，使得 \(\mathrm{H}_{d + j}^{(q + 1)} = 0\) 对所有整数 \(j\geqslant 0\) 成立

\begin{enumerate}
\def\labelenumi{\alph{enumi})}
\setcounter{enumi}{4}
\tightlist
\item
  设 \(k\) 是 \(\mathbf{K}\) 的一个子域，并把 \(\mathbf{K}(\mathbf{X})\) , \(k(\mathbf{X})\) 和 \(k[[\mathbf{X}]]\) 与 \(\mathbf{K}((\mathbf{X}))\) 的子环等同起来。证明 \(k[[\mathbf{X}]] \cap \mathbf{K}(\mathbf{X}) = k[[\mathbf{X}]] \cap k(\mathbf{X})\) 。
\end{enumerate}

\begin{enumerate}
\def\labelenumi{\arabic{enumi})}
\setcounter{enumi}{1}
\tightlist
\item
  设 \(a_1, a_2, \ldots, a_p\) 为整数 \(> 0\) ；用 \(\alpha_n\) 表示满足方程的有限序列 \((x_i)_{1 \leqslant i \leqslant p}\) （由整数 \(\geqslant 0\) 组成）的个数
\end{enumerate}

\[
a _ {1} x _ {1} + a _ {2} x _ {2} + \dots + a _ {p} x _ {p} = n.
\]

证明形式幂级数 \(\sum_{n=0}^{\infty} \alpha_n X^n\) （在 \(\mathbf{Q}\) 上）是有理分式

\[
\frac {1}{(1 - \mathrm{X} ^ {a _ {1}}) (1 - \mathrm{X} ^ {a _ {2}}) \dots (1 - \mathrm{X} ^ {a _ {p}})}.
\]

\begin{enumerate}
\def\labelenumi{\arabic{enumi})}
\setcounter{enumi}{2}
\tightlist
\item
  设 \(\mathbf{F}\) 是一个由整数 \(> 0\) 组成的有限集，并设 \(\beta_{n}\) 为有限序列 \((x_{i})\) （至多 \(n\) 项，所有项均属于 \(\mathbf{F}\) 且满足 \(\sum_{i} x_{i} = n\) ）的个数。证明形式幂级数 \(\sum_{n=0}^{\infty} \beta_n X^n\) 在 \(Q\) 上是有理分式
\end{enumerate}

\[
\frac {1}{1 - \sum_ {a \in F} X ^ {a}}.
\]

\begin{enumerate}
\def\labelenumi{\arabic{enumi})}
\setcounter{enumi}{3}
\item
  设 E 是特征为 2 的域 K 上具有无限基的向量空间；设 B 为这个空间上的外代数 \(\Lambda E\)，它是一个交换环。给出形式幂级数 \(u \in B[[X]]\) 的一个例子，使得 \(u^2 = 0\)，但不存在 B 的元素 \(\gamma \neq 0\) 使 \(\gamma u = 0\) 。
\item
  设 \(\mathbf{E} = \mathbf{A}[\mathbf{X}_1, \dots, \mathbf{X}_p]\) , \(\mathbf{F} = \mathbf{A}[[\mathbf{X}_1, \dots, \mathbf{X}_p]]\) .
\end{enumerate}

\begin{enumerate}
\def\labelenumi{\alph{enumi})}
\item
  如果 D 是从 E 到 F 的 Z-导子，则我们有 \(\omega(Du) \geqslant \omega(u) - 1\)，其中 \(u\) 是 E 中任意非零元。b) 如果 D' 是 F 的 Z-导子，则 \(\omega(D'v) \geqslant \omega(v') - 1\)，其中 \(v\) 是 F 中任意非零元。
\item
  如果赋予 E 由 F 诱导的拓扑，则 D 和 D' 是连续的。d) E 到 F 的每一个 Z-导子都以唯一的方式延拓为 F 的 Z-导子。e) 设 D 是 F 的 A-导子所成的 F-模。如果 D ∈ D 且 DX\_i = u\_i，则 D = ∑\_\{i=1\}\^{}\{p\} u\_i D\_i。
\end{enumerate}

\(f)\) \((\mathbf{D}_1,\dots ,\mathbf{D}_p)\) 是 F-模 \(\mathcal{D}\) 的一组基。

\begin{enumerate}
\def\labelenumi{\arabic{enumi})}
\setcounter{enumi}{5}
\tightlist
\item
  设 \((u_{\lambda})_{\lambda \in L}\) 是 \(\mathbf{A}[[\mathbf{X}_i]]_{i\in I}\) 的元素的一个族。如果族 \((1 + u_{\lambda})_{\lambda \in L}\) 是可乘的，则族 \((u_{\lambda})_{\lambda \in L}\) 是可和的。
\end{enumerate}

* 7) 若 A 是既约的，则 \(\mathbf{A}[[(\mathbf{X}_i)_{i\in I}]]\) 也是既约的。*

\begin{enumerate}
\def\labelenumi{\arabic{enumi})}
\setcounter{enumi}{7}
\tightlist
\item
  我们用 \((1 + \mathbf{X})^{\mathrm{T}}\) 表示元素 \(\exp (\mathrm{T}\log (1 + \mathbf{X}))\)，它属于 \(\mathbf{Q}[[\mathbf{X},\mathrm{T}]]\) 。
\end{enumerate}

\begin{enumerate}
\def\labelenumi{\alph{enumi})}
\tightlist
\item
  证明 \((1 + \mathbf{X})^{\mathrm{T}} = \sum_{n\geqslant 0}\frac{\mathrm{T}(\mathrm{T} - 1)\dots(\mathrm{T} - n + 1)}{n!}\mathbf{X}^{n}\) 。（\(\mathbf{X}^n\) 的系数是一个
\end{enumerate}

关于 T 的多项式，其值在 \(\mathbf{T} = n\in \mathbf{N}\) 处为已知。）

\begin{enumerate}
\def\labelenumi{\alph{enumi})}
\setcounter{enumi}{1}
\tightlist
\item
  证明 \((1 + \mathbf{X})^{\mathrm{T}}(1 + \mathbf{X})^{\mathrm{T}} = (1 + \mathbf{X})^{\mathrm{T} + \mathrm{T}}\) 。明确写出由此得到的二项式系数之间的恒等式。证明 \((1 + \mathbf{X} + \mathbf{Y} + \mathbf{XY})^{\mathrm{T}} = (1 + \mathbf{X})^{\mathrm{T}}(1 + \mathbf{Y})^{\mathrm{T}}\) 。
\end{enumerate}

